% s-167-p167-m002; original PDF 167
\Aphorism{s-167-p167-m002}{255}
Über sich selbst lachen können, gehört zur Menschlichkeit, mehr noch als über sich selbst Gericht halten können. Lachen: aber nicht, wie der „Idealist“ lacht, der mit bösem Auge auf sich hinschielt, um sich bei irgend einer Unvollkommenheit, einer Apostase oder Schauspielerei zu ertappen, sondern unschuldig und wohlwollend, als Einer, der gerade über seine Vollkommenheit und Idealität hinaus ist und aufgehört hat, sich in irgend einem für sich oder Andere einschränkenden Sinne wichtig zu nehmen. „Der Herr von Matthisson muss nicht denken, er wäre es, und ich muss nicht denken, ich wäre es“, sagt Goethe in seiner grundgütigen Art zu Eckermann. Wer irgendwie und irgendwann einmal dieser selbst-ironisirenden Gesellschaft von \Greek{θεῶν} \Greek{ῥεῖα} \Greek{ζῳόντων}, dieser goethisch-göttlichen Gesellschaft höchster, aber nicht pathetisch gestauter und gespannter Persönlichkeiten zugehört hat (auf die Dauer kann Niemand, auch sie selbst nicht, so leben und empfinden), der hält überhaupt keine andere Art von Ironie und Heiterkeit mehr aus, am wenigsten das sardonische Lachen, die krampfhafte, mit verzerrtem Munde grinsende Verhöhnung der menschlichen Ohnmacht in sich oder Anderen.
\par

